
Comme @@M0@@ est régulière, et @@M1@@ et @@M2@@ lisses, on en déduit d'après 1.3 que @@M3@@ et @@M4@@ sont régulières.

La proposition qui suit ne sera pas utilisée dans la suite du Séminaire :

\textbf{Proposition 1.4.} \uline{Soit} @@M5@@ \uline{un morphisme plat localement de présentation finie}. \uline{Pour que} @@M6@@ \uline{soit un morphisme d'intersection complète au sens de} 1.1, \uline{il faut et il suffit qu'il soit un morphisme d'intersection complète au sens de} EGA IV 19.3.6.

Supposons que @@M7@@ soit un morphisme d'intersection complète au sens de 1.1, et soit @@M8@@. Alors il existe un voisinage ouvert @@M9@@ de @@M10@@, et une factorisation :
@@M11@@
où @@M12@@ est une immersion régulière, et @@M13@@ un morphisme lisse. Soit @@M14@@ ; alors l'immersion @@M15@@ est régulière en @@M16@@. En effet, on peut supposer, quitte à restreindre @@M17@@, qu'il existe un @@M18@@-Module libre de type fini @@M19@@, et un homomorphisme régulier @@M20@@ tel que le complexe de Koszul associé à @@M21@@ soit une résolution de @@M22@@. On considère alors l'homomorphisme @@M23@@ induit sur la fibre en @@M24@@, et le complexe de Koszul associé @@M25@@. Comme @@M26@@ est lisse, donc plat, @@M27@@ est une résolution @@M28@@-plate de @@M29@@. Par suite, les groupes d'homologie de @@M30@@ sont les @@M31@@, où @@M32@@ est le corps résiduel de @@M33@@. Comme @@M34@@ est un morphisme plat, ils sont nuls, et @@M35@@ est régulière en @@M36@@. Par ailleurs, comme @@M37@@ est lisse, l'anneau local de @@M38@@ en @@M39@@ est régulier. Comme l'anneau local de @@M40@@ en @@M41@@ est quotient d'un anneau régulier par un idéal régulier,
